State some examples for aggregation rules in the modular approach for capital requirements

Let \( EC_i \) be the Economic Capital for module \( i, i = 1, \ldots, n \) and let EC denote the aggregated capital.

A risk aggregation rule is a mapping

\( f(EC_1, \ldots, EC_n) = EC \)

which maps the individual capital amounts per module to the total capital.

Aggregation Rules:
Simple Summation: \( EC = EC_1 + \ldots + EC_n \)
Correlation Adjusted Summation:
\( EC = \sqrt{\sum_{i=1}^{n} \sum_{j=1}^{n} \rho_{i,j} EC_i EC_j} \)

where \( \rho_{i,j} \in [0, 1] \) denote correlation parameters.
- This aggregation rule is also called Square Root Formula. 
  Risk aggregation in the Solvency II Standard Formula is based on the Square Root Formula.
- The Square Root Formula can be derived at the level of probability distributions in a multivariate Gaussian setting 
  or more generally for elliptical distributions.


State the Solvency II SCR formula and its relation to the Square Root formula 

Solvency Capital Requirement (SCR) in Solvency II, is defined by the Mean Value at Risk:

\( \text{SCR}(X) = \text{V@R}_{0.005}(X - \mathbb{E}[X]). \)

{Facts:
 The square root formula resembles the structure of the formula of the variance \(\text{Var}\) or standard deviation \(\sigma = \sqrt{\text{Var}}\), respectively, of two random variables with linear correlation \(\rho\):
\(
\text{Var}(X + Y) = \text{Var}(X) + \text{Var}(Y) + 2\rho \sqrt{\text{Var}(X)} \sqrt{\text{Var}(Y)},
\)
\(
\sigma(X + Y) = \sqrt{\sigma^2(X) + \sigma^2(Y) + 2\rho \sigma(X)\sigma(Y)}.
\)
\(X \sim \mathcal{N}(\mu, \sigma^2) \Rightarrow \text{V@R}_\lambda(X) = -\mu - \Phi^{-1}(\lambda)\sigma, \text{SCR}(X) = -\Phi^{-1}(\lambda)\sigma\)

 If \(X, Y\) are multivariate Gaussian risks with correlation \(\rho\), then
\(
\text{SCR}(X + Y) = \sqrt{(\text{SCR}(X))^2 + (\text{SCR}(Y))^2 + 2\rho\text{SCR}(X)\text{SCR}(Y)}.
\)
For risks satisfying a multivariate normal distribution (or more generally an elliptical distribution), the square root formula is correct for risk aggregation.
ATTENTION:
For other marginal distributions (in particular: skewed distributions, heavy tails) or dependencies 
(in particular: tail dependencies) the square root formula is not correct!


How can risk aggregation by fully probabilistically modelled, what are the drawbacks?

Framework: The total portfolio is given by \(X := X_1 + \ldots + X_n\).
Calculation of the total risk: margins-plus-copula approach:
Model the marginal distributions of \(X_1, \ldots, X_n\).
Model the dependency between \(X_1, \ldots, X_n\) in terms of a suitable copula. 
Simulate the distribution of \(X\) based on the marginal distributions and the copula.
Estimate the total risk capital EC in terms of a suitable risk measure.

This approach is quite challenging. In particular, the choice/specification of the copula may have a severe impact on the risk figures

Explain the basic idea of capital allocation

Mathematical description
Consider \(n\) risky sub-portfolios whose outcomes are given by \(X_1, \ldots, X_n\).
Then the total portfolio is given by \(X := X_1 + \cdots + X_n\).
If risk is measured in terms of a sub-additive risk measure \(\rho\), then there is a diversification effect in the total portfolio:
\(
\rho(X) \leq \sum_{i=1}^n \rho(X_i)
\)

Goal
Allocate to each sub-portfolio \(X_i\), \(i = 1, \ldots, n\), a risk capital contribution \(\rho(X_i \mid X)\) which reflects its contribution 
to the total risk \(\rho(X)\) and to the diversification effect.

Key properties of a meaningful capital allocation
Complete capital allocation: \(\sum_{i=1}^n \rho(X_i \mid X) = \rho(X)\)
Diversification: For each sub-portfolio \(X_i\), \(i = 1, \ldots, n\), the allocated capital is at most as large as the sub-portfolio’s standalone capital:
\(
\rho(X_i \mid X) \leq \rho(X_i) \quad \text{for } i = 1, \ldots, n.
\)

Give examples of capital allocation principles

Proportional allocation:
\(
\rho(X_i \mid X) = \frac{\rho(X_i)}{\sum_{j=1}^n \rho(X_j)} \rho(X)
\)

Marginal principles, e.g.:
Euler allocation (for homogeneous risk measures):
\(
\rho(X_i \mid X) = \left. \frac{d}{dh} \rho(X + h X_i) \right|_{h=0}
\)
 Covariance principle:
\(
\rho(X_i \mid X) = \frac{\text{Cov}(X_i, X)}{\sigma^2(X)} \rho(X)
\)

Game theoretical methods, e.g., Shapley allocation


Explain the details of the Euler Principal for capital allocation
and also what useful properties it satisfies for risk measures


Euler's rule for positive homogeneous functions:
Let \( f : \mathbb{R}^n \setminus \{0\} \to \mathbb{R} \) be differentiable and positive homogeneous of order \(\alpha > 0\), i.e.,
\(
f(t\mathbf{w}) = t^\alpha f(\mathbf{w}) \quad \text{for all } t > 0, \mathbf{w} \in \mathbb{R}^n \setminus \{0\}.
\)
Differentiating on both sides of this equation with respect to \(t\) yields
\(
\alpha t^{\alpha-1} f(\mathbf{w}) = \nabla f(t\mathbf{w}) \cdot \mathbf{w} = \sum_{i=1}^n \frac{\partial f}{\partial w_i} (t\mathbf{w}) \cdot w_i
\)

hence for \( t = 1 \)
\(\alpha f(\mathbf{w}) = \nabla f(\mathbf{w}) \cdot \mathbf{w} = \sum_{i=1}^n \frac{\partial f}{\partial w_i} (\mathbf{w}) \cdot w_i\).

Application to positive homogeneous risk measures (\(\alpha = 1\)):
Let \(\mathcal{O} \subset \mathbb{R}^n \setminus \{0\}\) be an open set of portfolio weights with \(\mathbf{1} \in \mathcal{O}\), 
and let \(\rho\) be a positive homogeneous risk measure.
Define the weighted portfolio \(\mathbf{X}(\mathbf{w}) = \sum_{i=1}^n w_i X_i, \mathbf{w} \in \mathcal{O}\), and the corresponding risk functional
\(
f_\rho(\mathbf{w}) = \rho(\mathbf{X}(\mathbf{w})).
\)
Since \(f_\rho\) is positive homogeneous of order 1, Euler's rule implies
\(
f_\rho(\mathbf{w}) = \sum_{i=1}^n \frac{\partial f_\rho}{\partial w_i} (\mathbf{w}) \cdot w_i, .
\)
hence
 \(\rho(\mathbf{X}) = f_\rho(1) = \sum_{i=1}^n \frac{\partial f_\rho}{\partial w_i} (1)\)

Let \( f_\rho \) be positive homogeneous and differentiable at \mathbf{w} = 1. Then the allocation

\[
\rho(X_i \vert X) := \frac{\partial f_\rho}{\partial w_i}(1), \quad i = 1, \ldots, n,
\]

is called the Euler Capital Allocation Principle.

Remarks:
1. The Euler principle satisfies property 1 (complete capital allocation).
2. Property 2 (diversification) is satisfied if \rho is convex.
3. The Euler principle plays also a prominent role for the characterization of RoRac compatibility.


State 2 important Risk-Adjusted Performance Measures

Sharpe ratio:

- The (ex ante) Sharpe Ratio of an investment with random one-period return \( R \) is defined by
\(
SR(R) = \frac{\mathbb{E}[R] - r_0}{\sigma(R)}.
\)

The Sharpe Ratio is a measure which relates the risk premium \(\mathbb{E}[R] - r_0\), defined as the excess of the average return over the risk-free interest rate, 
to the corresponding volatility \(\sigma(R)\).

Return on Risk adjusted Capital:

The RoRaC (ex ante) the ratio between the expected economic profit and the corresponding risk capital:
\(
\text{RoRaC} := \frac{\text{expected economic profit}}{\text{Risk Adjusted Capital}}.
\)

What is a RoRaC compatible capital allocation principle

Definition: For a given risk measure \(\rho\), let

\[
\text{RoRaC}(X) = \frac{\mathbb{E}[X]}{\rho(X)} 
\quad \text{and} \quad 
\text{RoRaC}(X_k \vert X) = \frac{\mathbb{E}[X_k]}{\rho(X_k \vert X)}
\]

denote the RoRaC of the total portfolio \(X\) and of the sub-portfolio \(X_k\) respectively. A capital allocation \(\rho(X_k \vert X)\) is said to be RoRaC compatible if there exists an \(\varepsilon > 0\) such that

\[
\text{RoRaC}(X_k \vert X) > \text{RoRaC}(X) 
\Rightarrow 
\text{RoRaC}(X + hX_k) > \text{RoRaC}(X) 
\quad \text{for all } 0 < h < \varepsilon.
\]

Interpretation:

- Typically, \(X\) signifies a quantity with \(\rho(X) > 0\) such as the net asset value at time \(1\) minus initial capital at time \(0\). In other words, \(\rho(X)\) could, for example, be the solvency capital requirement.
- Increasing the amount of sub-portfolio \(X_k\), whose RoRaC exceeds RoRaC of the total portfolio, increases the RoRaC of the total portfolio.
